\documentclass[11pt,a4paper]{article}
\usepackage[margin=2.5cm]{geometry}
\usepackage[utf8]{inputenc}
\usepackage[T1]{fontenc}
\usepackage{lmodern}
\usepackage{booktabs}
\usepackage{longtable}
\usepackage{array}
\usepackage{hyperref}
\usepackage{xcolor}
\usepackage{listings}

\lstset{
  basicstyle=\ttfamily\small,
  breaklines=true,
  frame=single,
  columns=fullflexible,
  keepspaces=true,
  showstringspaces=false,
}

\title{Password Manager (POO)\\[0.3em]\large Secure Single-User CLI Vault}
\author{Academic exercise in object-oriented programming and applied cryptography}
\date{}

\begin{document}
\maketitle
\tableofcontents
\bigskip

\noindent
A local, offline, single-user password manager written in \textbf{Java 21}. It keeps
all credentials in one authenticated, encrypted file (\texttt{boveda.json}) and is
operated entirely from the terminal. There is no server, no cloud account and no
telemetry: confidentiality and integrity depend only on your master password and
standard cryptography.

The project demonstrates three security indicators end-to-end:

\begin{center}
\begin{tabular}{@{}lll@{}}
\toprule
Indicator & Primitive \\
\midrule
IND1 --- Key encapsulation     & RSA-3072 OAEP-SHA256 wrapping a random AES-256 DEK \\
IND2 --- Digital signatures   & Ed25519 over the canonical vault payload \\
IND3 --- Zero-knowledge proof & Interactive Schnorr proof over RFC 3526 Group 14 \\
\bottomrule
\end{tabular}
\end{center}

\section{What this project does}

In plain terms:

\begin{enumerate}
  \item \textbf{You create a vault once.} On first launch you choose a username and a
        master password. The program generates all cryptographic material and writes
        \texttt{boveda.json}.
  \item \textbf{You unlock the vault with your master password.} The program verifies the
        file's signature first, then proves you know the master password using a
        zero-knowledge proof. A wrong password is rejected; the password itself is
        never stored.
  \item \textbf{You manage credentials.} From a text menu you can add a credential
        (service, username, notes, password), list entries without revealing
        passwords, or view one entry's full details.
  \item \textbf{You save and exit.} Saving re-encrypts the credential list, signs the
        payload and writes the file atomically; exiting wipes the session key and
        decrypted credentials from memory.
  \item \textbf{You can change the master password.} Only the small private-key
        material is re-encrypted and the data key re-wrapped; the (potentially large)
        credentials blob is \textbf{not} re-encrypted.
\end{enumerate}

Everything is stored locally in a single JSON file. Losing the file means losing the
vault; losing the master password means the vault cannot be recovered.

\subsection{Feature summary}

\begin{longtable}{@{}p{5.2cm}p{9cm}@{}}
\toprule
Feature & Behaviour \\
\midrule
First-run registration & Username + master password (with confirmation) creates a new vault \\
Signature-first login & Ed25519 signature verified \textbf{before} the password is requested \\
Password authentication & Three-move Schnorr proof; no password hash stored \\
Add credential & Service, username, notes, optional generated password \\
List credentials & Shows service, username and notes; never the password \\
View credential & Full details for one selected entry, including password \\
Password generator & \texttt{SecureRandom}, 15 characters plus \texttt{-} separators \\
Master-password rotation & Re-encrypts private keys, re-wraps DEK, keeps credentials blob \\
Secure input & \texttt{System.console().readPassword()} with \texttt{Scanner} fallback \\
Atomic persistence & Temporary file + atomic rename, owner-only permissions on POSIX \\
Memory hygiene & Secrets as \texttt{char[]}/\texttt{byte[]}, wiped on lock/exit where possible \\
Standalone evidence & Three executable classes for IND1, IND2, IND3 \\
Automated tests & 63 JUnit~5 tests, including abuse scenarios \\
\bottomrule
\end{longtable}

\section{Quick start}

\subsection{Prerequisites}
\begin{itemize}
  \item JDK \textbf{21} or newer
  \item Maven \textbf{3.9+}
\end{itemize}

\subsection{Run the tests}
\begin{lstlisting}
mvn clean test
\end{lstlisting}
Expected: \textbf{63 tests, 0 failures, 0 errors}.

\subsection{Run the application}
\begin{lstlisting}
mvn exec:java
\end{lstlisting}
\texttt{cli.Main} is the default main class. To use a specific vault path:
\begin{lstlisting}
mvn exec:java -Dexec.args="/path/to/boveda.json"
\end{lstlisting}
On first run the program asks for a username and a master password (twice) and
creates \texttt{boveda.json} in the current directory (or at the given path). On
later runs it asks only for the master password and opens the menu.

\subsection{Build a jar}
\begin{lstlisting}
mvn clean package
\end{lstlisting}
The resulting jar is \textbf{not} a fat jar; dependencies must be on the classpath.
The manifest main class is \texttt{cli.Main}.

\paragraph{Secure input note.} \texttt{System.console().readPassword()} masks input
with a real terminal console. Without a console (IDE run configurations, pipes,
Maven \texttt{exec:java}) the program falls back to line-based \texttt{Scanner} and
input may be echoed. For daily use, run from a terminal or an IDE console.

\section{How to use it}

\subsection{First run (registration)}
\begin{lstlisting}
No vault found. First-run registration.
Username: alice
Master password: ********
Confirm master password: ********
Vault created: /home/alice/boveda.json
\end{lstlisting}
The program generates a random salt, derives the master key with Argon2id, computes
the Schnorr verifier, creates RSA and Ed25519 key pairs, encrypts the private keys,
generates a random data key, signs the payload and writes the vault.

\subsection{Later runs (login)}
\begin{lstlisting}
Master password: ********
Vault unlocked.
\end{lstlisting}
The signature is checked before the password prompt. A wrong password is rejected;
after three failed attempts the program exits with a security error.

\subsection{Main menu}
\begin{lstlisting}
=== Password Manager ===
1. Add credential
2. List credentials
3. View credential details
4. Change master password
5. Save, wipe memory and exit
0. Exit without saving
Option:
\end{lstlisting}

\begin{description}
  \item[Option 1 --- Add credential] Prompts: \texttt{Service}, \texttt{Username/email},
    \texttt{Notes}, \texttt{Generate a password? (y/N)}, \texttt{Password}. If you
    answer \texttt{y}, a password is generated, printed for copying, and stored
    without asking you to type one. Each added credential is saved immediately.
  \item[Option 2 --- List credentials] Shows \texttt{n. service | username | notes};
    stored passwords are never printed.
  \item[Option 3 --- View credential] Select an entry by number to print service,
    username, notes and password. The password is copied into a temporary
    \texttt{char[]}, printed, then wiped.
  \item[Option 4 --- Change master password] Prompts for current, new and confirmed
    new master password. Current password is checked with ZKP; private keys are
    re-encrypted under a new KEK from a fresh salt; verifier recomputed; DEK
    re-wrapped; payload re-signed; encrypted credentials reused byte-for-byte; old
    password stops working immediately.
  \item[Option 5 / 0 --- Exit] \texttt{5} saves, wipes session key, exits;
    \texttt{0} exits without saving (read-only session).
\end{description}

\section{How it works}

\subsection{Session lifecycle}
\begin{lstlisting}
lock --register/login--> unlocked --menu operations--> save --> locked
                          |
                          +-- exit / error --> wipe DEK, keys, credentials
\end{lstlisting}
While locked, no DEK, private key or decrypted credential is kept. While unlocked,
the DEK and decrypted data live in process memory (trusted-RAM window).
\texttt{Gestor.lock()} zeroizes the DEK buffer, clears credential passwords and
drops private-key references.

\subsection{Data flow}
\begin{lstlisting}
master password (char[])
        |
        +-- Argon2id(password, salt, 64 bytes)
        |       +-- bytes  0..31  -> KEK (AES-256)
        |       +-- bytes 32..63  -> Schnorr secret seed
        |                            +- x in [1, q-1]  ->  Y = g^x mod p
        |
        +-- KEK --AES-256-GCM--> RSA private key
        |                     +-> Ed25519 private key
        |
        +-- random DEK (32 bytes) <--RSA-OAEP-SHA256-- KEM public key
                 |
                 +--AES-256-GCM--> credentials JSON blob
\end{lstlisting}

\subsection{Registration}
\begin{enumerate}
  \item Generate a 16-byte random salt.
  \item Derive 64 bytes with Argon2id.
  \item First 32 bytes = KEK; last 32 bytes = Schnorr seed.
  \item Compute verifier $Y = g^x \bmod p$ and store it (public).
  \item Generate RSA-3072 KEM key pair and Ed25519 signing key pair.
  \item Encrypt both private keys with AES-256-GCM under the KEK.
  \item Generate random 32-byte DEK; wrap it with RSA-OAEP under KEM public key.
  \item Sign canonical payload with Ed25519; write vault atomically.
\end{enumerate}

\subsection{Login}
\begin{enumerate}
  \item Read \texttt{boveda.json}.
  \item Verify Ed25519 signature \textbf{before} asking for the password; mismatch
        raises \texttt{SecurityException} immediately.
  \item Prompt for master password.
  \item Run three-move Schnorr proof against stored verifier: prover sends
        commitment $R = g^r \bmod p$; verifier sends random challenge $c$; prover
        sends $s = (r + c \cdot x) \bmod q$; verifier checks
        $g^s \equiv R \cdot Y^c \pmod p$.
  \item Only on success: derive KEK, decrypt private keys, decapsulate DEK, decrypt
        credential blob and keep it in memory.
\end{enumerate}

\subsection{Saving}
\begin{enumerate}
  \item Serialize credentials to JSON (passwords as numeric UTF-16 code units,
        never strings).
  \item Encrypt JSON with AES-256-GCM under DEK (fresh 96-bit nonce).
  \item Wrap DEK with RSA-OAEP under KEM public key.
  \item Build payload; sign canonical JSON with Ed25519.
  \item Write to temporary file; atomically move over \texttt{boveda.json}.
\end{enumerate}

\subsection{Master-password rotation}
\begin{enumerate}
  \item Re-derive KEK and Schnorr seed from new password with fresh salt.
  \item Re-encrypt existing RSA and Ed25519 private keys under new KEK.
  \item Recompute verifier $Y$.
  \item Re-wrap the same DEK under unchanged KEM public key.
  \item Re-sign payload.
  \item Keep existing \texttt{credentialsIv} and \texttt{credentialsCiphertext} unchanged.
\end{enumerate}

\section{Architecture}

\subsection{Packages and responsibilities}
\begin{longtable}{@{}p{5.5cm}p{8.7cm}@{}}
\toprule
Package / class & Responsibility \\
\midrule
\texttt{cli.Main} & Interactive menu, secure input, registration/login, rotation \\
\texttt{cli.PasswordManagerApp} & Registers the Bouncy Castle provider \\
\texttt{crypto.Argon2KDF} & Argon2id key derivation (64 MiB, 3 iters, 16-byte salt) \\
\texttt{crypto.VaultCipher} & AES-256-GCM encryption/decryption with fresh nonces \\
\texttt{crypto.KeyEncapsulator} & RSA-3072 OAEP-SHA256 DEK encapsulation (IND1) \\
\texttt{crypto.DigitalSigner} & Ed25519 key generation, signing, verification (IND2) \\
\texttt{crypto.SchnorrZKP} & Schnorr registration and interactive proof (IND3) \\
\texttt{model.Credential} & Credential POJO with \texttt{char[]} password \\
\texttt{model.Usuario} & Base user type holding public identity data \\
\texttt{model.UsuarioMaestro} & Salt, verifier, public keys, KEK-encrypted private keys, rotation \\
\texttt{model.Gestor} & Singleton session manager: unlock, save, lock, rotation \\
\texttt{storage.VaultFile} & Gson DTO for on-disk payload and signature \\
\texttt{storage.PersistenciaBoveda} & Save/load, signature verification, decryption, re-wrap \\
\texttt{util.SecretWiper} & Best-effort zeroization of \texttt{byte[]}, \texttt{char[]}, keys \\
\texttt{util.GeneradorDeContrase\~nas} & \texttt{SecureRandom} password generator \\
\texttt{evidencias.*} & Standalone IND1/IND2/IND3 demonstration programs \\
\bottomrule
\end{longtable}

\subsection{Key hierarchy}
\begin{itemize}
  \item Master password $\to$ Argon2id $\to$ \textbf{KEK} + Schnorr seed.
  \item KEK $\to$ AES-256-GCM $\to$ RSA and Ed25519 private keys.
  \item Random DEK $\to$ AES-256-GCM $\to$ credentials blob.
  \item DEK $\to$ RSA-OAEP-SHA256 $\to$ \texttt{encapsulatedDek}.
  \item Schnorr seed $\to x \to Y = g^x \bmod p$ (public verifier).
\end{itemize}
The DEK is independent of the master password, which makes rotation cheap: the
password-derived material changes, but the DEK and credentials blob do not.

\section{Vault file format}
\texttt{boveda.json} is UTF-8 JSON (abridged):
\begin{lstlisting}
{
  "payload": {
    "usuario": {
      "salt": "<base64 16 bytes>",
      "verifier": "<base64 Y>",
      "kemPublicKey": "<base64 X.509 RSA>",
      "kemPrivateKeyIv": "<base64 12 bytes>",
      "kemPrivateKeyCiphertext": "<base64 AES-GCM>",
      "signingPublicKey": "<base64 X.509 Ed25519>",
      "signingPrivateKeyIv": "<base64 12 bytes>",
      "signingPrivateKeyCiphertext": "<base64 AES-GCM>",
      "nombreDeUsuario": "alice"
    },
    "encapsulatedDek": "<base64 RSA-OAEP>",
    "credentialsIv": "<base64 12 bytes>",
    "credentialsCiphertext": "<base64 AES-GCM>"
  },
  "signature": "<base64 Ed25519 over the canonical payload>"
}
\end{lstlisting}
Credential passwords are serialized as numeric UTF-16 code units inside the
encrypted blob --- never as JSON strings. The file is written atomically
(temporary file + rename) with owner-only permissions on POSIX.

\section{Security design}

\subsection{Protected assets}
\begin{longtable}{@{}p{4.6cm}p{3.6cm}p{6cm}@{}}
\toprule
Asset & Where it lives & Protection \\
\midrule
Master password & User memory / terminal only & Never stored; derives KEK + ZKP secret \\
Credential passwords & \texttt{credentialsCiphertext} & AES-256-GCM under random DEK \\
Credential metadata & Same encrypted blob & AES-256-GCM \\
KEM private key (RSA) & \texttt{kemPrivateKeyCiphertext} & AES-256-GCM under KEK \\
Signing private key (Ed25519) & \texttt{signingPrivateKeyCiphertext} & AES-256-GCM under KEK \\
Data-encryption key (DEK) & \texttt{encapsulatedDek} & RSA-OAEP-SHA256 under KEM public key \\
Vault integrity & \texttt{signature} & Ed25519 over canonical payload \\
ZKP verifier $Y$ & \texttt{usuario.verifier} & Public; reveals nothing about password \\
\bottomrule
\end{longtable}

\subsection{Threat model and trust assumptions}
\begin{longtable}{@{}p{3.2cm}p{4.4cm}p{6.6cm}@{}}
\toprule
Adversary & Capability & Mitigation \\
\midrule
Offline disk attacker & Read/copy/modify/replace/delete vault & Argon2id slows guessing; Ed25519 rejects modification; AES-GCM rejects tampering; RSA-OAEP protects DEK \\
Password guesser & Has vault; tries candidate passwords offline & Argon2id (64 MiB, 3 iters); Schnorr verifier is not a fast hash oracle \\
File tamperer & Flips bytes in signature/ciphertext/metadata & Signature verified before decryption; mismatch raises \texttt{SecurityException} \\
Impersonator & Runs Schnorr without master password & Wrong-secret proofs fail verification \\
Wrong-key attacker & Unwraps DEK with unrelated RSA key & RSA-OAEP unwrap fails \\
\bottomrule
\end{longtable}

\paragraph{Explicit trust assumptions.}
RAM during an active session is trusted (memory scraping, debuggers, core dumps,
swap, hibernation out of scope; \texttt{lock()} wipes best-effort). OS, JVM,
\texttt{SecureRandom}, Bouncy Castle and Gson are trusted. Terminal is trusted
(\texttt{readPassword()} avoids echo; keyloggers/shoulder surfing not prevented).
Master password has reasonable entropy. Machine is not compromised while vault is
open.

\subsection{Cryptographic choices}
\begin{longtable}{@{}p{3.6cm}p{3.4cm}p{3.8cm}p{3.4cm}@{}}
\toprule
Purpose & Primitive & Parameters & Why \\
\midrule
Password KDF & Argon2id (RFC 9106) & 64 MiB, 3 iters, parallelism 1, 16-byte salt, 64-byte output & Memory-hard, side-channel resistant; expensive offline guessing \\
Vault encryption & AES-256-GCM & 256-bit key, 96-bit random nonce, 128-bit tag & AEAD confidentiality + integrity; no padding oracle \\
Vault authentication & Ed25519 & 32-byte public key, 64-byte signature & Deterministic, fast, small, misuse-resistant \\
DEK wrapping (IND1) & RSA-3072 OAEP & SHA-256 + MGF1-SHA256 & CCA-secure encapsulation with explicit hash pinning \\
Password proof (IND3) & Schnorr ZKP & RFC 3526 Group 14: 2048-bit $p$, $q=(p-1)/2$, $g=2$ & Proves knowledge of $x$ without revealing password or hash \\
Randomness & \texttt{SecureRandom} & JCA default & Salts, nonces, keys, challenges \\
\bottomrule
\end{longtable}

\subsection{Known limitations}
\begin{itemize}
  \item Vault substitution / key substitution: Ed25519 public key stored inside
        signed payload; whole-file replacement not detected without key pinning /
        out-of-band fingerprint (TOFU is future work).
  \item Rollback and deletion: older valid vault can be restored; no monotonic counter.
  \item Metadata leakage: file size / ciphertext length reveal approximate count/size.
  \item Memory wiping is best-effort; \texttt{BigInteger} secrets cannot be zeroized.
  \item Single user, local only; no multi-user access control, sync, recovery, audit
        log, GUI or clipboard integration.
  \item Not production-ready: teaching/academic implementation.
\end{itemize}

\section{Micro-evidence programs}
Three independent \texttt{main} classes under \texttt{src/main/java/evidencias/}
demonstrate IND1, IND2 and IND3 in isolation, printing inputs, outputs and explicit
\texttt{PASS}/\texttt{FAIL} lines, ending with \texttt{[RESULT] PASS}.
\begin{lstlisting}
mvn exec:java -Dexec.mainClass=evidencias.EvidenciaKEM
mvn exec:java -Dexec.mainClass=evidencias.EvidenciaFirma
mvn exec:java -Dexec.mainClass=evidencias.EvidenciaZKP
\end{lstlisting}
Output is reproducible: fixed demo DEK/password/salt and the RFC 8032 Ed25519 test
vector; randomized values are not printed. IND1 checks recovered DEK equals input,
tampered encapsulation rejected, unrelated key rejected. IND2 checks RFC 8032
vector, valid signature accepted, one-byte tampering rejected, different message
rejected. IND3 checks verifier in order-$q$ subgroup, deterministic registration,
valid proof accepted, impersonation rejected, tampered response rejected.

\section{Automated tests}
\texttt{mvn clean test} runs \textbf{63 JUnit~5 tests}:
\begin{itemize}
  \item \textbf{crypto}: Argon2id vectors/parameters, AES-GCM round-trip and
        tamper/wrong-key rejection, RSA-OAEP KEM, Ed25519 sign/verify, Schnorr
        proof and impersonation rejection.
  \item \textbf{model}: \texttt{Credential} secret handling, \texttt{UsuarioMaestro}
        registration/ZKP/unlock, \texttt{Gestor} session lifecycle, DEK zeroization,
        master-password rotation.
  \item \textbf{storage}: save/reload round-trip, tampered-signature abort,
        verify-before-decrypt, wrong-DEK AEAD failure.
  \item \textbf{security}: full end-to-end lifecycle + abuse scenarios: one-byte file
        tampering, one-byte AES-GCM tampering (\texttt{AEADBadTagException}), ZKP
        impersonation, unrelated KEM private key.
  \item \textbf{cli}: scripted registration/login/menu sessions across restarts,
        wrong-password lockout, password rotation through the menu.
  \item \textbf{util}: secure password generator.
\end{itemize}
Notable negative tests: flip one byte of Ed25519 signature $\to$ \texttt{SecurityException}
before decryption; flip one byte of AES-GCM ciphertext $\to$ \texttt{AEADBadTagException};
wrong password $\to$ proof verification returns \texttt{false}; decapsulate DEK with
unrelated RSA key $\to$ \texttt{GeneralSecurityException}.

\section{Project structure}
\begin{lstlisting}
src/main/java/
  cli/          Main (interactive CLI), PasswordManagerApp (provider bootstrap)
  crypto/       Argon2KDF, VaultCipher, KeyEncapsulator, DigitalSigner, SchnorrZKP
  model/        Credential, Usuario, UsuarioMaestro, Gestor
  storage/      VaultFile (DTO), PersistenciaBoveda (save/load/verify/rewrap)
  util/         SecretWiper, GeneradorDeContraseñas
  evidencias/   EvidenciaKEM, EvidenciaFirma, EvidenciaZKP
src/test/java/  crypto, model, storage, security, cli, util
pom.xml
README.md
\end{lstlisting}

\section{Licensing and legal considerations}
\begin{tabular}{@{}ll@{}}
\toprule
Component & License \\
\midrule
Bouncy Castle (\texttt{bcprov-jdk18on}) & MIT-style permissive license \\
Gson (\texttt{com.google.code.gson:gson}) & Apache License 2.0 \\
JUnit Jupiter & Eclipse Public License 2.0 \\
Apache Maven plugins & Apache License 2.0 \\
\bottomrule
\end{tabular}

\bigskip
If redistributing the project or a binary bundling dependencies: keep Bouncy
Castle copyright/license notice; keep Gson \texttt{LICENSE}/\texttt{NOTICE} per
Apache~2.0; add a \texttt{LICENSE} file (academic coursework; MIT is a compatible
simple choice). Standard public cryptography is used; verify local export/import
regulations before distributing binaries.

\section{Project status}
\begin{tabular}{@{}cll@{}}
\toprule
Phase & Scope & State \\
\midrule
1 & Maven/Java 21 scaffolding, Bouncy Castle provider & Done \\
2 & Cryptographic primitives and unit tests & Done \\
3 & Domain model and authenticated JSON persistence & Done \\
4 & Interactive CLI, registration/login, credential management, rotation & Done \\
5 & End-to-end and abuse-scenario tests & Done \\
6 & Standalone evidence programs and technical documentation & Done \\
\bottomrule
\end{tabular}

\bigskip
\textbf{63 tests pass} (\texttt{mvn clean test}). The three evidence programs run
cleanly and produce reproducible output. Intended for academic demonstration, not
production deployment.

\end{document}
